\section{গণিত Mathematics}
লিখিত সংখ্যা: 12 এবং ১২।
\[
 x^2+12=14,\qquad x^{২}+১২=১৪,\qquad x_{১২}^{৩_{৪}}
\]
% Single-token scripts and arguments must work just like Western digits.
\[
 x^২+x_৩+২^৩+\frac১২+\sqrt৪
\]
\begingroup
  \setbox0=\hbox{$x^২+x_৩+২^৩+\frac১২+\sqrt৪$}
  \setbox2=\hbox{$x^{২}+x_{৩}+২^{৩}+\frac{১}{২}+\sqrt{৪}$}
  \ifdim\wd0=\wd2\typeout{MATH-SINGLE-TOKEN-WIDTH: equal}\fi
\endgroup
\[
 \frac{১২}{৩}+\sqrt{৪}+\sqrt[৩]{৮}
 \qquad \begin{pmatrix}১&২\\৩&৪\end{pmatrix}
\]
\[
 \text{বেগ}=\frac{\text{দূরত্ব}}{\text{সময়}},\qquad
 A_{\text{কাঠুরিয়া}},\qquad x_{y_{\text{বাংলা}}}
\]
\begin{align}
 x&=১২+3\label{eq:test}\\
 \intertext{এখানে বাংলা এবং English আছে।}
 y&=\begin{cases}১&\text{যদি }x>০\\০&\text{অন্যথায়}\end{cases}
\end{align}
সমীকরণ \eqref{eq:test}।
\textbf{গাঢ় লেখা এবং $x^{২}+১২$।}
